%% 本文件由 工程文件/scripts/assemble.py 生成，勿手改。

%% begin units/ch09_u01.tex
% ch09_u01 —— 原书 p.370--p.374（第 9 章全部：9.1.1 -- 9.1.13）
% 原书章标题「9. To be inserted into Section 6」：本章是待插入第 6 节的材料。

\bdchap{ch:ins6}{待插入第 6 节的材料}

\SourcePage{370}

\bdsec{9.1. 定理 8.1.4 与命题 8.1.5 的证明}

沿用 5.2.13、8.1.1 与 8.1.4 的记号。定理 8.1.4 与命题 8.1.5 容易由下述陈述推出。

\bdp{9.1.1} \emph{定理.} $L_0$ 在 $H^{i}(U_{\chi},I_{\chi}\lambda_{\mathcal{L}}^{-1})$ 上的特征值 $\ge -d(\chi)/2$。若 $i>0$，则这些特征值 $>-d(\chi)/2$。若 $i=0$，则特征值 $-d(\chi)/2$ 的重数为 1，相应的特征空间是 (392) 的像。

下面开始证明该定理。记 $I^{U}_{\chi}$ 为 $I_{\chi}$ 到 $U_{\chi}$ 的限制，即 $I^{U}_{\chi}$ 是右 $\mathcal{D}$-模 $\omega_{\operatorname{Orb}_{\chi}}$ 沿闭嵌入 $\operatorname{Orb}_{\chi}\hookrightarrow U_{\chi}$ 的正向像。考虑 $I^{U}_{\chi}\lambda_{\mathcal{L}}^{-1}$ 上的 $\mathcal{O}$-模滤链，其第 $k$ 项由支撑在 $\operatorname{Orb}_{\chi}$ 的第 $k$ 个无穷小邻域上的截面构成。该滤链是 $\operatorname{Aut}^{0}_{2}O$-不变的，且 $\operatorname{gr}_{j}(I^{U}_{\chi}\lambda_{\mathcal{L}}^{-1})=\omega_{\operatorname{Orb}_{\chi}}\otimes\lambda_{\mathcal{L}}^{-1}\otimes\operatorname{Sym}^{j}\mathcal{N}_{\chi}$，其中 $\mathcal{N}_{\chi}$ 是 $\operatorname{Orb}_{\chi}\subset U_{\chi}$ 的法层。利用 (241) 得到 $\operatorname{Aut}^{0}_{2}O$-等变同构 $\operatorname{gr}_{j}(I^{U}_{\chi}\lambda_{\mathcal{L}}^{-1})=\mathfrak{d}_{\mathcal{L},\chi}\otimes\operatorname{Sym}^{j}\mathcal{N}_{\chi}$。由 4.6.15，$L_0$ 在 $\mathfrak{d}_{\mathcal{L},\chi}$ 上乘以 $-d(\chi)/2$。于是只剩证明下述结论。
% [Tutor Note: 原书 p.370 此段两处印作 Aut^0_2 O（上标 0、外加下标 2），而 9.1.2 的注与 9.1.5 等处印作 Aut^0 O；原文未说明二者关系，此处照录底本，未作统一。]

\bdp{9.1.2} \emph{命题.}
\begin{enumerate}[label=\roman*),leftmargin=2.2em]
\item $L_0$ 在 $H^{i}(\operatorname{Orb}_{\chi},\operatorname{Sym}^{j}\mathcal{N}_{\chi})$ 上的特征值非负。
\item 若 $i>0$ 或 $j>0$，则它们为正。$\operatorname{Orb}_{\chi}$ 上除常数外没有 $L_0$-不变的正则函数。
\end{enumerate}

\emph{注.} $L_0$ 在 $H^{i}(\operatorname{Orb}_{\chi},\operatorname{Sym}^{j}\mathcal{N}_{\chi})$ 上的特征值是整数，因为 $\mathcal{N}_{\chi}$ 是 $\operatorname{Aut}^{0}O$-等变层。

证明该命题之前，先给出几个引理。

\bdp{9.1.3} 先引入一些记号。回忆 $\chi$ 是 $G$ 的支配余权。固定嘉当子群 $H\subset G$ 与包含 $H$ 的 Borel 子群 $B\subset G$。以下把「余权」理解为「$H$ 的余权」，把「支配」理解为「关于 $B$ 支配」。记 $t^{\chi}\in H(K)$ 为 $t\in\mathbb{C}((t))^{*}=K^{*}$ 在 $\chi:\mathbb{G}_{m}\to H$ 下的像。回忆 $\operatorname{Orb}_{\chi}$ 是\SourcePage{371}$[\chi]$ 的 $G(O)$-轨道，其中 $[\chi]$ 是 $t^{\chi}$ 在 $\mathcal{GR}=G(K)/G(O)$ 中的像。

记 $\operatorname{orb}_{\chi}$ 为 $[\chi]$ 的 $G$-轨道，$P^{-}_{\chi}$ 为 $[\chi]$ 在 $G$ 中的稳定化子，即 $P^{-}_{\chi}=\{g\in G\mid t^{-\chi}gt^{\chi}\in G(O)\}$。$P^{-}_{\chi}$ 是 $G$ 的抛物子群，其 $\operatorname{Lie}P^{-}_{\chi}$ 是 $\operatorname{Lie}H$ 与满足 $(\alpha,\chi)\le 0$ 的根 $\alpha$ 所对应的根空间之和（特别地，$P^{-}_{\chi}$ 包含与 $B$ 相对的 Borel 子群 $B^{-}\supset H$）。于是 $\operatorname{orb}_{\chi}=G/P^{-}_{\chi}$ 是射影簇。显然 $\operatorname{Aut}^{0}O$ 在 $\operatorname{orb}_{\chi}$ 上的作用平凡。

\bdp{9.1.4} $O$ 的自同态构成仿射半群概形 $\operatorname{End}^{0}O$（对 $\mathbb{C}$-代数 $R$，$\operatorname{End}^{0}O$ 的 $R$-点是一个 $R$-态射 $f:R[[t]]\to R[[t]]$，满足 $f(t)\in tR[[t]]$）。$\operatorname{Aut}^{0}O$ 在 $\operatorname{End}^{0}O$ 中稠密。记 $\mathbf{0}\in\operatorname{End}^{0}O$ 为 $O=\mathbb{C}[[t]]$ 的自同态 $t\mapsto 0$。

\bdp{9.1.5} \emph{引理.}
\begin{enumerate}[label=\roman*),leftmargin=2.2em]
\item $\operatorname{Aut}^{0}O$ 在 $\operatorname{Orb}_{\chi}$ 上的作用可延拓为 $\operatorname{End}^{0}O$ 在 $\operatorname{Orb}_{\chi}$ 上的作用。
\item 设 $\varphi$ 是 $\operatorname{Orb}_{\chi}$ 上对应于 $\mathbf{0}\in\operatorname{End}^{0}O$ 的自同态。则 $\varphi^{2}=\varphi$，且 $\varphi$ 的不动点概形等于 $\operatorname{orb}_{\chi}$。
\item 由 $\varphi$ 诱导的态射 $p:\operatorname{Orb}_{\chi}\to\operatorname{orb}_{\chi}$ 是仿射的。其纤维同构于仿射空间。
\end{enumerate}

\emph{证明.}
\begin{enumerate}[label=\roman*),leftmargin=2.2em]
\item $\operatorname{Orb}_{\chi}=G(O)/S$，其中 $S$ 是 $[\chi]$ 在 $G(O)$ 中的稳定化子。$\operatorname{Aut}^{0}O$ 在 $G(O)$ 上的作用可延拓为 $\operatorname{End}^{0}O$ 的作用。因 $S$ 是 $\operatorname{Aut}^{0}O$-不变的，它也是 $\operatorname{End}^{0}O$-不变的。
\item 对应于 $\mathbf{0}\in\operatorname{End}^{0}O$ 的态射 $f:G(O)\to G(O)$ 是复合 $G(O)\to G\hookrightarrow G(O)$。故 $\varphi(\operatorname{Orb}_{\chi})\subset\operatorname{orb}_{\chi}$。显然 $\varphi$ 在 $\operatorname{orb}_{\chi}$ 上的限制等于 $\operatorname{id}$。
\item $G(O)=G\cdot U$，其中 $U:=\operatorname{Ker}(G(O)\to G)$。有 $f(S)\subset S$，故 $S=S_{G}\cdot S_{U}$，$S_{G}:=S\cap G$，$S_{U}:=S\cap U$。$p$ 是自然态射 $G(O)/S\to G(O)/(S_{G}\cdot U)=G/S_{G}=\operatorname{orb}_{\chi}$。因 $U$ 是投射幂幺的，$(S_{G}\cdot U)/S=U/S_{U}$ 同构于仿射空间。\hfill$\Box$
\end{enumerate}

\bdp{9.1.6} \emph{注.} 由 9.1.5(ii) 可知，$L_0$ 在 $\operatorname{Orb}_{\chi}$ 上的不动点概形等于 $\operatorname{orb}_{\chi}$。

\SourcePage{372}

\bdp{9.1.7} 因 $p:\operatorname{Orb}_{\chi}\to\operatorname{orb}_{\chi}$ 是仿射的，
\[
H^{i}(\operatorname{Orb}_{\chi},\operatorname{Sym}^{j}\mathcal{N}_{\chi})=H^{i}(\operatorname{orb}_{\chi},p_{*}\operatorname{Sym}^{j}\mathcal{N}_{\chi}).
\]
$p$ 是 $\operatorname{Aut}^{0}O$-等变的，故 $\operatorname{Aut}^{0}O$、从而 $L_0$ 在 $p_{*}\operatorname{Sym}^{j}\mathcal{N}_{\chi}$ 上有作用。为证明命题 9.1.2，只需证明下述引理。

\bdp{9.1.8} \emph{引理.} $L_0$ 在 $p_{*}\operatorname{Sym}^{j}\mathcal{N}_{\chi}$ 上的特征值非负。若 $j>0$，则它们为正。若 $j=0$，则 $L_0$ 的零特征层等于 $\operatorname{orb}_{\chi}$ 的结构层。

\emph{证明.} 记 $\mathcal{O}_{\operatorname{Orb}}$ 与 $\mathcal{O}_{\operatorname{orb}}$ 分别为 $\operatorname{Orb}_{\chi}$ 与 $\operatorname{orb}_{\chi}$ 的结构层。由 9.1.5(i) 知，$L_0$ 在 $p_{*}\mathcal{O}_{\operatorname{Orb}}$ 上的特征值非负。9.1.5(ii) 或 9.1.6 表明，$L_0:p_{*}\mathcal{O}_{\operatorname{Orb}}\to p_{*}\mathcal{O}_{\operatorname{Orb}}$ 的余核等于 $\mathcal{O}_{\operatorname{orb}}$。

显然的态射 $\mathcal{O}_{\operatorname{Orb}}\otimes(\mathfrak{g}\otimes K/\mathfrak{g}\otimes O)\to\mathcal{N}_{\chi}$ 是满射，且是 $\operatorname{Aut}^{0}O$-等变的。它诱导 $\operatorname{Aut}^{0}O$-等变满射 $p_{*}\mathcal{O}_{\operatorname{Orb}}\otimes\operatorname{Sym}^{j}(\mathfrak{g}\otimes(K/O))\to p_{*}\operatorname{Sym}^{j}\mathcal{N}_{\chi}$。因 $L_0$ 在 $K/O$ 上的特征值为正，引理得证。\hfill$\Box$

\bdp{9.1.9} 至此已证明 9.1.2，从而证明了 8.1.4 与 8.1.5。下面用态射 $p:\operatorname{Orb}_{\chi}\to\operatorname{orb}_{\chi}$ 计算 $\operatorname{Orb}_{\chi}$ 的典范丛。答案（见 9.1.12、9.1.13）将在 10.1.7 中使用。

\bdp{9.1.10} $\operatorname{Orb}_{\chi}$ 是 $G(O)$ 的齐性空间，而 $\operatorname{orb}_{\chi}$ 是 $G$ 的齐性空间。利用投影 $G(O)\to G(O/tO)=G$，得到 $G(O)$ 在 $\operatorname{orb}_{\chi}$ 上的作用。态射 $p:\operatorname{Orb}_{\chi}\to\operatorname{orb}_{\chi}$ 是 $G(O)$-等变的。\footnote{当然，嵌入 $\operatorname{orb}_{\chi}\hookrightarrow\operatorname{Orb}_{\chi}$ 不是 $G(O)$-等变的。是否需要这条脚注？}
% [Tutor Note: 原书 p.372 脚注末尾印有作者手稿中的批注「DO WE NEED THIS FOOTNOTE?」，底本如此，保留原样。]

\bdp{9.1.11} \emph{命题.} 函子 $p^{*}$ 诱导 $\operatorname{orb}_{\chi}$ 上 $G$-等变线丛的群胚与 $\operatorname{Orb}_{\chi}$ 上 $G(O)$-等变线丛的群胚之间的等价。

\SourcePage{373}

\emph{证明.} 有 $\operatorname{Orb}_{\chi}=G(O)/S$，$\operatorname{orb}_{\chi}=G/S_{G}$，其中 $S$ 是 $[\chi]$ 在 $G(O)$ 中的稳定化子，$S_{G}=S\cap G$。事实上，$S_{G}$ 是 $S$ 在 $G$ 中的像，而 $p:G(O)/S\to G/S_{G}$ 由投影 $G(O)\to G$ 诱导（见 9.1.5(iii) 的证明）。须证明态射 $\pi:S\to S_{G}$ 诱导同构 $\operatorname{Hom}(S_{G},\mathbb{G}_{m})\to\operatorname{Hom}(S,\mathbb{G}_{m})$。这是显然的，因为 $\operatorname{Ker}\pi\subset\operatorname{Ker}(G(O)\to G)$ 是投射幂幺的。\hfill$\Box$

\emph{注.} 之所以对等变丛陈述该命题，是因为将按这种形式使用它。当然，若去掉「等变」一词，陈述仍然成立（事实上，$p$ 是局部平凡纤维化，其纤维同构于仿射空间）。此外，若 $G$ 单连通，则 $\operatorname{orb}_{\chi}$ 上的线丛有唯一的 $G$-等变结构（因为由 9.1.3，$\operatorname{orb}_{\chi}=G/P^{-}_{\chi}$，而 $P^{-}_{\chi}$ 是抛物的）。

\bdp{9.1.12} 典范层 $\omega_{\operatorname{Orb}_{\chi}}$ 是 $\operatorname{Orb}_{\chi}$ 上的 $G(O)$-等变线丛。由 9.1.11，它来自 $\operatorname{orb}_{\chi}$ 上唯一的 $G$-等变线丛 $\mathcal{M}_{\chi}$。因 $\operatorname{orb}_{\chi}=G/P^{-}_{\chi}$（见 9.1.3），$\operatorname{orb}_{\chi}$ 上 $G$-等变线丛的同构类由 $\operatorname{Hom}(P^{-}_{\chi},\mathbb{G}_{m})$ 参数化。嵌入 $H\hookrightarrow P^{-}_{\chi}$ 诱导嵌入 $\operatorname{Hom}(P^{-}_{\chi},\mathbb{G}_{m})\hookrightarrow\operatorname{Hom}(H,\mathbb{G}_{m})$。于是 $\mathcal{M}_{\chi}$ 定义 $H$ 的一个权，它可看作元素 $l_{\chi}\in\mathfrak{h}^{*}$。

\bdp{9.1.13} \emph{命题.} $l_{\chi}=B\chi$，其中 $\chi\in\operatorname{Hom}(\mathbb{G}_{m},H)$ 按通常方式与 $\mathfrak{h}$ 的元素等同，$B:\mathfrak{h}\to\mathfrak{h}^{*}$ 是对应于内积 (18) 的线性算子。

\emph{证明.} $\operatorname{Orb}_{\chi}$ 在 $[\chi]$ 处的切空间等于
\begin{equation}\tag{393}
(\mathfrak{g}\otimes O)/((\mathfrak{g}\otimes O)\cap t^{\chi}(\mathfrak{g}\otimes O)t^{-\chi}).
\end{equation}
$H$ 在 (393) 上的作用来自 $H$ 在 $\mathfrak{g}\otimes O$ 上的伴随作用。故 (393) 中出现的 $H$ 的权是正根，且对正根 $\alpha$，它在 (393) 中的重数等于 $(\chi,\alpha)$。因此，与 (393) 的对偶向量空间的行列式相对应的 $\mathfrak{h}$ 的权等于\SourcePage{374}
\[
-\sum_{\alpha>0}(\chi,\alpha)\cdot\alpha=-\frac{1}{2}\sum_{\alpha}(\chi,\alpha)\cdot\alpha=B\chi.
\]
\hfill$\Box$

给作者的注记：记法 $U:=\operatorname{Ker}(G(O)\to G)$ 与记法 $U_{\chi}$ 并不完全相容。这样可以吗？
% [Tutor Note: 原书 p.374 此处印有作者写给合作者的批注「Note for the authors: ... Is this OK ???」，底本如此，保留原样。]
%% end units/ch09_u01.tex
